\section{Literature Review and Study of Existing Systems}
\label{ch:lit}

\subsection{Study of Existing Systems}
\textbf{Offline messengers.} \emph{Bridgefy}, a Bluetooth messenger used at large protests, is Study~11 of Table~\ref{tab:studies}; \emph{FireChat} offered phone-to-phone chat and was withdrawn; \emph{Briar} needs no server and exchanges encrypted messages over Bluetooth or Wi-Fi when devices meet~\cite{briar}. \textbf{Dedicated radios.} \emph{goTenna Mesh} pairs phones with radio units; \emph{Meshtastic} uses LoRa with a controlled flood in which a node stays silent if a neighbour already repeats the packet~\cite{meshtastic}. \textbf{Research and humanitarian work.} The \emph{Serval Project} builds disaster phone meshes with its \emph{Rhizome} store-and-forward layer~\cite{serval}; the \emph{delay-tolerant networking (DTN)} architecture keeps data and carries it across intermittent links~\cite{fall2003dtn,rfc4838}. ``n/d'' marks a property the sources do not document.
\subsection{Review of Research Papers and Technical Resources}
\label{sec:papers}
Table~\ref{tab:studies} summarises the twelve sources we studied, what each contributes and what it means for EcoMesh.
\begingroup\footnotesize
\setlength{\tabcolsep}{3pt}
\begin{longtable}{@{}p{0.5cm}>{\raggedright\arraybackslash}p{3.5cm}p{6.3cm}p{6.0cm}@{}}
\caption{Twelve studies: technique, contribution and relevance to EcoMesh}\label{tab:studies}\\
\toprule \textbf{\#} & \textbf{Source} & \textbf{Technique and contribution} & \textbf{Relevance and limitation}\\ \midrule \endfirsthead
\toprule \textbf{\#} & \textbf{Source} & \textbf{Technique and contribution} & \textbf{Relevance and limitation}\\ \midrule \endhead
1 & DSDV, Perkins and Bhagwat 1994~\cite{perkins1994dsdv} & Each node keeps a table of distances in advance; numbered updates stop routing loops & Routes exist before the first SOS; counts hops only, and update traffic rises as nodes move\\
2 & DSR, RFC 4728~\cite{rfc4728dsr} & Routes are found on demand and the whole path is written into each packet & A carried path allows loop checks and the ``relayed via $n$'' display; nothing is held\\
3 & AODV, RFC 3561~\cite{rfc3561aodv} & A flooded request finds a destination and the reply sets up the route backwards & Counts hops; no holding of messages and no priority\\
4 & OLSR, RFC 3626~\cite{rfc3626olsr} & Link state maintained in advance; selected relays reduce repeated broadcasts & Nearest relative of EcoMesh; still uses hop count by default\\
5 & OSPF v2, RFC 2328~\cite{rfc2328ospf} & Numbered link-state flooding with Dijkstra run on each router & EcoMesh keeps this structure but uses a measured cost; designed for wired links\\
6 & Babel, RFC 8966~\cite{rfc8966babel} & Loop-free distance vector for wireless meshes with any metric plugged in & Shows the wireless metric must be replaceable; covers routing only\\
7 & ETX, De Couto et al. 2003~\cite{decouto2003etx} & Path cost from measured delivery ratios; shows that fewest-hop paths choose lossy links & Main basis for our cost; we add a penalty for acknowledgement failures\\
8 & Epidemic routing, Vahdat and Becker 2000~\cite{vahdat2000epidemic} & Copy every message to every node encountered & Delivers well but fills buffers; we use one bounded queue instead\\
9 & DTN architecture, Fall 2003; RFC 4838~\cite{fall2003dtn,rfc4838} & Nodes keep, carry and later pass on data, with custody handed over & Basis of our third exit; ignores priority and link quality\\
10 & Disaster-response networks, Reina et al. 2015~\cite{reina2015disaster} & Survey of ad hoc, vehicular, delay-tolerant and sensor networks in disasters & Supports our framing and the call for systems that are evaluated\\
11 & Breaking Bridgefy, Albrecht et al. 2021~\cite{albrecht2021bridgefy} & Attacks on a real mesh messenger: users tracked, messages unauthenticated, network shut down & Defines our threat model; motivates encryption and signatures\\
12 & AES-GCM, Ed25519, Dijkstra~\cite{nist80038d,bernstein2012ed25519,dijkstra1959} & Authenticated encryption, fast secure signatures and shortest paths & The primitives and routing algorithm we use\\
\bottomrule
\end{longtable}
\endgroup

\subsection{Comparative Analysis}
\label{sec:compare}
Table~\ref{tab:compare1} sets the systems of Section~2.1 and the protocols of Table~\ref{tab:studies} against EcoMesh: the infrastructure each needs, how it routes, whether it prioritises, whether it keeps a message that has no route, and how it protects content.
\begin{table}[!htbp]
\caption[Existing systems and routing protocols compared with EcoMesh]{Existing offline-communication systems and routing protocols compared with EcoMesh (n/d = not documented in the sources consulted)}
\label{tab:compare1}
\footnotesize
\begin{tabularx}{\textwidth}{@{}>{\raggedright\arraybackslash}p{2.1cm}>{\raggedright\arraybackslash}X>{\raggedright\arraybackslash}X>{\raggedright\arraybackslash}p{1.7cm}>{\raggedright\arraybackslash}p{1.9cm}>{\raggedright\arraybackslash}X@{}}
\toprule
\textbf{System} & \textbf{Infrastructure} & \textbf{Routing} & \textbf{Priority} & \textbf{Store-fwd} & \textbf{Encryption}\\ \midrule
Bridgefy & Bluetooth LE phones & n/d & n/d & n/d & none effective in analysed version~\cite{albrecht2021bridgefy}\\
Briar & none (Bluetooth/Wi-Fi, Tor) & direct device sync & n/d & on contact & end-to-end~\cite{briar}\\
Meshtastic & LoRa radios & managed flood~\cite{meshtastic} & n/d & n/d & channel-level\\
Serval & Wi-Fi ad hoc phones & B.A.T.M.A.N., later own & n/d & yes (Rhizome)~\cite{serval} & n/d\\
DTN & none & architecture only & bundle class of service & yes (custody) & separate specification\\
AODV, DSR, DSDV, OLSR, OSPF, Babel & none & hop count, configured or pluggable metric & none & none or short & not in protocol\\
\textbf{EcoMesh} & \textbf{none} & \textbf{link state, weighted Dijkstra} & \textbf{3 levels, every hop} & \textbf{yes, bounded} & \textbf{AES-GCM + Ed25519}\\
\bottomrule
\end{tabularx}
\end{table}

\subsection{Limitations of Existing Systems}
Hop count, the usual default metric, picks lossy links~\cite{decouto2003etx}; no system prioritises an emergency at each hop; ad hoc routing drops a message with no route, while delay-tolerant designs keep it but ignore link quality and priority; and deployed messengers publish little measurement, one with no effective confidentiality or authenticity~\cite{albrecht2021bridgefy}.
\subsection{Identified Research or Engineering Gap}
\label{sec:gap}
\textbf{None of the surveyed systems joins priority-aware, link-quality-weighted routing with a delay-tolerant fallback, and none is evaluated quantitatively under controlled failures.}
\subsection{Proposed Improvements}
EcoMesh routes over the cost $1+5(1-q)+3f$ (Section~\ref{sec:cost}), re-applies SOS priority at each hop, holds unroutable messages in a bounded queue that never evicts an SOS for routine traffic, uses vetted cryptography, and measures all of it, including weighted against hop-count routing.
